\exercisesection{\S~13}

\begin{enumerate}
\def\labelenumi{\arabic{enumi})}
\tightlist
\item
  可分裂单李代数的 \(\leq 80\) 的维数为：
\end{enumerate}

3 (A \(_{1}\) = B \(_{1}\) = C \(_{1}\) ), 8 (A \(_{2}\) ), 10 (B \(_{2}\) = C \(_{2}\) ), 14 (G \(_{2}\) ), 15 (A \(_{3}\) = D \(_{3}\) ), 21 (B \(_{3}\) 与 C \(_{3}\) ), 24 (A \(_{4}\) ), 28 (D \(_{4}\) ), 35 (A \(_{5}\) ), 36 (B \(_{4}\) 与 C \(_{4}\) ), 45 (D \(_{5}\) ), 48 (A \(_{6}\) ), 52 (F \(_{4}\) ), 55 (B \(_{5}\) 与 C \(_{5}\) ), 63 (A \(_{7}\) ), 66 (D \(_{6}\) ), 78 (B \(_{6}\) , C \(_{6}\) 与 E \(_{6}\) ), 80 (A \(_{8}\) ).

\begin{enumerate}
\def\labelenumi{\arabic{enumi})}
\setcounter{enumi}{1}
\tightlist
\item
  令 \((\mathfrak{g},\mathfrak{h})\) 为分裂单李代数，B 为 \(R(\mathfrak{g},\mathfrak{h})\) 的基。维数最小的单 \(\mathfrak{g}\)-模 \(E(\lambda)\)（\(\lambda \in P_{++} - \{0\}\)）是那些 \(\lambda\) 为下列权之一的模：
\end{enumerate}

\(\varpi_{1}\)（A \(_{1}\)）；\(\varpi_{1}\) 与 \(\varpi_{l}\)（A \(_{l}\)，\(l \geq 2\)）；\(\varpi_{1}\)（B \(_{l}\) 与 C \(_{l}\)，\(l \geq 2\)）；\(\varpi_{1}\)、\(\varpi_{3}\) 与 \(\varpi_{4}\)（D \(_{4}\)）；\(\varpi_{1}\)（D \(_{l}\)，\(l \geq 5\)）；\(\varpi_{1}\) 与 \(\varpi_{6}\)（E \(_{6}\)）；\(\varpi_{7}\)（E \(_{7}\)）；\(\varpi_{8}\)（E \(_{8}\)）；\(\varpi_{4}\)（F \(_{4}\)）；\(\varpi_{1}\)（G \(_{2}\)）。任何两个这样的模都可经 g 的一个自同构相互变换。

\(E_{8}\) 型是唯一使伴随表示达到最小维数的类型。

\begin{enumerate}
\def\labelenumi{\arabic{enumi})}
\setcounter{enumi}{2}
\tightlist
\item
  \begin{enumerate}
  \def\labelenumii{\alph{enumii})}
  \tightlist
  \item
    定义一个从 \(\mathfrak{sl}(4,k)\) 到正交李代数 \(\mathfrak{o}_S(6,k)\) 的同构。（利用第1节.V 的表示 \(\bigwedge^2\sigma\) 是正交的且维数为 6 这一事实。）\(\mathfrak{sl}(4,k)\) 的两个次数为 4 的不可约表示类型对应于 \(\mathfrak{o}_S(6,k)\) 的两个半旋量表示。
  \end{enumerate}
\end{enumerate}

\begin{enumerate}
\def\labelenumi{\alph{enumi})}
\setcounter{enumi}{1}
\item
  定义一个从 \(\mathfrak{sp}(4,k)\) 到正交李代数 \(\mathfrak{o}_S(5,k)\) 的同构。（利用第3节.V 的表示 \(\sigma_2\) 是正交的且维数为 5 这一事实。）\(\mathfrak{sp}(4,k)\) 的次数为 4 的不可约表示对应于 \(\mathfrak{o}_S(5,k)\) 的旋量表示。
\item
  定义一个同构 \(\mathfrak{sl}(2,k)\times\mathfrak{sl}(2,k)\to\mathfrak{o}_{S}(4,k)\)（利用两个 \(\mathfrak{sl}(2,k)\) 因子的恒等表示的张量积）。借助第I章§6习题26重新得到这一结果。
\end{enumerate}

\begin{enumerate}
\def\labelenumi{\arabic{enumi})}
\setcounter{enumi}{3}
\tightlist
\item
  设 \(S\) 为 \(n\) 阶方阵
\end{enumerate}

\[
(\delta_ {i, n + 1 - i}) = \left( \begin{array}{c c c c c} 0 & 0 & \dots & 0 & 1 \\ 0 & 0 & \dots & 1 & 0 \\ \vdots & \vdots & \ddots & \vdots & \vdots \\ 0 & 1 & \dots & 0 & 0 \\ 1 & 0 & \dots & 0 & 0 \end{array} \right).
\]

\(\mathfrak{o}_{S}(n,k)\) 的元素是关于副对角线反对称的矩阵 \((a_{ij})\)：

\[
a _ {i, j} = - a _ {n + 1 - j, n + 1 - i} \quad \text {   for   every   pair   } (i, j).
\]

李代数 \(\mathfrak{o}_S(n,k)\) 是 \(\mathrm{D}_{n/2}\) 型的可分裂单李代数（若 \(n\) 为偶数且 \(\geq 6\)），而是 \(\mathrm{B}_{(n-1)/2}\) 型（若 \(n\) 为奇数且 \(\geq 5\)）。\(\mathfrak{o}_S(n,k)\) 的对角（分别地，上三角）元素构成分裂嘉当（分别地，Borel）子代数。

\begin{enumerate}
\def\labelenumi{\arabic{enumi})}
\setcounter{enumi}{4}
\item
  记号沿用第1节.IV，\(\mathrm{A}_l\) 型。证明：若 \(n\) \(\geq 0\)，则 \(\mathbf{S}^n\sigma\) 是 \(\mathfrak{sl}(l + 1,k)\) 的最高权为 \(n\varpi_1\) 的不可约表示。（把 \(\mathbf{S}^n\sigma\) 实现在次数为 \(n\) 的齐次多项式（关于 \(\mathrm{X}_0,\dots ,\mathrm{X}_l\)）的空间上，并注意在相差一个位似的意义下，唯一的多项式 \(f\)（满足 \(\partial f / \partial \mathrm{X}_i = 0\)，对 \(i\geq 1\)）是 \(\mathrm{X}_0^n\)。）证明这个表示的所有权都是重数 1 的。
\item
  记号沿用第2节.IV，\(\mathrm{B}_l\) 型。证明：若 \(1 \leq r \leq l - 1\)，则 \(\operatorname{E}(\varpi_r)\) 的维数是 \(\binom{2l + 1}{r}\)，并由此给出 \(\bigwedge^r \sigma\) 是最高权为 \(\varpi_r\) 的基本表示这一事实的另一个证明。
\item
  记号沿用第2节.VII，\(\mathrm{B}_l\) 型。证明 \(\mathbf{O}_0^+ (\Psi)\) 是 \(\mathbf{SO}(\Psi)\) 的换位子群。（注意 \(\mathbf{O}(\Psi)\) 等于 \(\{\pm 1\} \times \mathbf{SO}(\Psi)\)，因而与 \(\mathbf{SO}(\Psi)\) 有相同的换位子群，然后应用《代数学》第IX章§9习题11 b）。）由此推出 \(\operatorname{Aut}_e(\mathfrak{g}) = \mathbf{O}_0^+ (\Psi)\)。
\item
  记号沿用第3节.IV，\(C_{l}\) 型（\(l \geq 1\)）。证明 \(S^{2}\sigma\) 等价于 g 的伴随表示。
\item
  记号沿用第3节.VII，\(C_l\) 型（\(l \geq 1\)）。特别地，我们把 \(\mathrm{Aut}_e(\mathfrak{g})\) 与 \(\mathbf{Sp}(\varPsi)/\{\pm 1\}\) 的一个子群等同起来。证明辛平延（《代数学》第IX章§4习题6）在 \(\mathbf{Sp}(\varPsi)/\{\pm 1\}\) 中的像属于 \(\mathrm{Aut}_e(\mathfrak{g})\)。由此推出 \(\mathrm{Aut}_e(\mathfrak{g}) = \mathbf{Sp}(\varPsi)/\{\pm 1\}\)（《代数学》第IX章§5习题11），并且 \(\mathrm{Aut}(\mathfrak{g})/\mathrm{Aut}_e(\mathfrak{g})\) 可以与 \(k^*/k^{*2}\) 等同起来。
\end{enumerate}

¶ 10) 记号沿用第4节.IV，D \(_{l}\) 型（l ≥ 2）。

\begin{enumerate}
\def\labelenumi{\alph{enumi})}
\tightlist
\item
  令 \(x\) 与 \(y\) 为 \(\bigwedge^{l} \mathrm{V}\) 中由下式定义的元素
\end{enumerate}

\[
x = e _ {1} \wedge \dots \wedge e _ {l - 1} \wedge e _ {l} \quad \text { and } \quad y = e _ {1} \wedge \dots \wedge e _ {l - 1} \wedge e _ {- l}.
\]

元素 x 是权为 \(2\varpi_{l}\) 的本原元，y 是权为 \(2\varpi_{l-1}\) 的本原元。\(\bigwedge^{l}V\) 中由 x（分别地，y）生成的子模 X（分别地，Y）同构于 \(\mathrm{E}(2\varpi_{l})\)（分别地，\(\mathrm{E}(2\varpi_{l-1})\)）。通过计算维数证明 \(\bigwedge^{l}V = X \oplus Y\)；特别地，\(\bigwedge^{l}V\) 是两个互不同构的单模之和。

\begin{enumerate}
\def\labelenumi{\alph{enumi})}
\setcounter{enumi}{1}
\tightlist
\item
  令 \(e = e_{1} \wedge \cdots \wedge e_{l} \wedge e_{-1} \wedge \cdots \wedge e_{-l} \in \bigwedge^{2l}V\)，并令 \(\varPsi_{l}\) 为 \(\varPsi\) 到 \(\bigwedge^{l}V\) 上的扩张。设 \(z \in \bigwedge^{l}V\)。证明下列等价关系：
\end{enumerate}

\[
z \in X \Longleftrightarrow z \wedge t = \Psi_ {l} (z, t) e \quad \text {   for   all   } t \in \bigwedge^ {l} V
\]

\[
z \in \mathrm{Y} \Longleftrightarrow z \wedge t = - \Psi_ {l} (z, t) e \quad \text {   for   all   } t \in \bigwedge^ {l} \mathrm{V}.
\]

（若以 \(X'\) 与 \(Y'\) 表示由右边定义的子空间，先证明 \(X'\) 与 \(Y'\) 在 \(\mathfrak{g}\) 作用下稳定，且分别含有 \(x\) 与 \(y\)。）

\begin{enumerate}
\def\labelenumi{\alph{enumi})}
\setcounter{enumi}{2}
\item
  假设 z 是纯元素（《代数学》第III章§11第13节），并以 \(M_{z}\) 表示与之相伴的 V 的 l 维子空间。证明 \(M_{z}\) 全迷向当且仅当 z 属于 X 或 Y。（利用下述事实：当 \(M_{z}\) 全迷向时，存在一个正交变换把 \(M_{x}\) 变到它；当 \(M_{z}\) 不全迷向时，构造一个 l-向量 t 使得 \(z \wedge t = 0, \Psi_{l}(z, t) = 1\)，并应用上面的 b）。）当 \(z \in X\)（分别地，\(z \in Y\)）时，\(\mathrm{M}_{x}/(\mathrm{M}_{x} \cap \mathrm{M}_{z})\) 的维数是偶（分别地，奇）整数，参见《代数学》第IX章§6习题18 d）。
\item
  设 s 是 V 的一个正向（分别地，反向）相似变换。证明 \(\bigwedge^{l}s\) 保持 X 与 Y 稳定（分别地，互换 X 与 Y）。
\end{enumerate}

\begin{enumerate}
\def\labelenumi{\arabic{enumi})}
\setcounter{enumi}{10}
\item
  记号沿用第4节.VII，\(\mathrm{D}_l\) 型（\(l \geq 3\)）。证明 \(\mathbf{O}_0^+(\varPsi)\) 是 \(\mathbf{SO}(\varPsi)\) 的换位子群。（应用《代数学》第IX章§6习题17 b）及《代数学》第IX章§9习题11 b）。）由此推出 \(\mathrm{Aut}_e(\mathfrak{g})\) 等于 \(\mathbf{O}_0^+(\varPsi)\) 在 \(\mathbf{SO}(\varPsi)/\{\pm 1\}\) 中的像。此外，\(-1\) 属于 \(\mathbf{O}_0^+(\varPsi)\) 当且仅当 \(l\) 为偶数或 \(-1\) 是 \(k\) 中的平方（《代数学》第IX章§9习题11 c））。
\item
  \(\mathfrak{sl}(n,k)\) 的 Killing 型是 \((X,Y)\mapsto 2n\operatorname {Tr}(XY)\)。\(\mathfrak{sp}(n,k)\)（\(n\) 为偶数时）的 Killing 型是 \((X,Y)\mapsto (n + 2)\operatorname {Tr}(XY)\)。\(\mathfrak{o}_S(n,k)\)（其中 \(S\) 是秩为 \(n\) 的非退化对称矩阵）的 Killing 型是 \((X,Y)\mapsto (n - 2)\operatorname {Tr}(XY)\)。
\item
  \(\mathfrak{g}\) 的不变多项式函数代数由下列元素生成：
\end{enumerate}

\(a\) ) 在 \(\mathrm{A}_l\) 型情形，由函数 \(X\mapsto \mathrm{Tr}(X^i)\)（\(2\leq i\leq l + 1\)）生成；

\begin{enumerate}
\def\labelenumi{\alph{enumi})}
\setcounter{enumi}{1}
\item
  在 \(B_{l}\) 型情形，由函数 \(X \mapsto \operatorname{Tr}(X^{2i})\)（\(1 \leq i \leq l\)）生成；
\item
  在 \(C_{l}\) 型情形，由函数 \(X \mapsto \operatorname{Tr}(X^{2i})\)（\(1 \leq i \leq l\)）生成；
\end{enumerate}

\(d\) ) 在 \(\mathrm{D}_l\) 型情形，由函数 \(X \mapsto \operatorname{Tr}(X^{2i})\)（\(1 \leq i \leq l - 1\)）以及两个多项式函数 \(\tilde{f}\) 之一（满足 \(\tilde{f}(X)^2 = (-1)^l \det(X)\)）生成。

\begin{enumerate}
\def\labelenumi{\arabic{enumi})}
\setcounter{enumi}{13}
\tightlist
\item
  \begin{enumerate}
  \def\labelenumii{\alph{enumii})}
  \tightlist
  \item
    令 G 为按 §7 习题26 的程序与 \(\mathfrak{g} = \mathfrak{sl}(n, k)\) 相伴的群。\(k^{n}\) 上的自然 g-模结构给出一个同态 \(\varphi : G \to \mathbf{GL}(n, k)\)。利用该处 h) 证明 \(\varphi\) 是单射，并利用该处 f) 证明
  \end{enumerate}
\end{enumerate}

\[
\operatorname{Im} (\varphi) = \mathbf {S L} (n, k).
\]

\begin{enumerate}
\def\labelenumi{\alph{enumi})}
\setcounter{enumi}{1}
\item
  令 E 为有限维 \(\mathfrak{sl}(n,k)\)-模，\(\rho\) 为 \(\mathfrak{sl}(n,k)\) 的相应表示。证明存在唯一的表示 \(\pi:\mathbf{SL}(n,k)\to\mathbf{GL}(\mathrm{E})\)，使得 \(\pi(e^{x})=e^{\rho(x)}\) 对 \(\mathfrak{sl}(n,k)\) 的每个幂零元 x 成立（利用 a)）。我们称 \(\rho\) 与 \(\pi\) 相容。把 §1 第4节中对 n=2 证明的结果加以推广。
\item
  假设 \(k\) 是 \(\mathbf{R}, \mathbf{C}\)，或是一个关于离散赋值完备且剩余域特征 \(\neq 0\) 的域。证明 \(\rho\) 与 \(\pi\) 相容当且仅当 \(\pi\) 是李群同态且 \(L(\pi) = \rho\)（方法同 §1 习题18 b)）。
\end{enumerate}

\(d\) ) 对 \(\mathfrak{sp}(2n,k)\) 与 \(\mathbf{Sp}(2n,k)\) 证明类似的结果。

¶15) 令 V 为有限维 \(\geq 2\) 的向量空间，g 为 End(V) 的李子代数，\(\theta\) 为 g 的元素。我们作如下假设：

\begin{enumerate}
\def\labelenumi{(\roman{enumi})}
\item
  V 是半单 g-模；
\item
  \(\theta\) 的秩为 1（即 \(\dim \operatorname{Im}(\theta) = 1\)）；
\item
  直线 \(\operatorname{Im}(\theta)\) 生成 U(g)-模 V。
\end{enumerate}

\begin{enumerate}
\def\labelenumi{\alph{enumi})}
\item
  证明这些假设成立（对适当选取的 \(\theta\)）：当 \(\mathfrak{g} = \mathfrak{sl}(V)\)、当 \(\mathfrak{g} = \mathfrak{gl}(V)\)，或当 V 上存在非退化交错双线性形式 \(\varPsi\) 使得 \(\mathfrak{g} = \mathfrak{sp}(\varPsi)\) 或 \(\mathfrak{g} = k.1 \oplus \mathfrak{sp}(\varPsi)\) 时。在每种情形，都可取 \(\theta\) 为幂零元；仅在第二种情形（且仅在该情形）可取 \(\theta\) 为半单元。
\item
  我们将证明上述四种情形是仅有的可能情形。立即可化为 k 代数闭的情形。证明此时 V 是单 g-模，且 \(g = c \oplus s\)，其中 s 半单，c = 0 或 c = k.1；证明 V 不同构于维数 \(\geq 2\) 的 g-模的张量积；由此推出 s 是单的。
\item
  取定 \(\mathfrak{h}\) 为 \(\mathfrak{s}\) 的一个嘉当子代数，以及 \(\mathrm{R}(\mathfrak{s},\mathfrak{h})\) 的基 B。令 \(\lambda\) 为 \(\mathfrak{s}\)-模 V 的（关于 B 的）最高权，并令 \(e\) 为 V 中权为 \(\lambda\) 的非零元素。对偶模 \(\mathrm{V}^*\) 的最高权是 \(\lambda^{*} = -w_{0}\lambda\)（§7 第5节）；令 \(e^*\) 为 \(\mathrm{V}^*\) 中权为 \(\lambda^*\) 的非零元素。按惯例把 \(\mathrm{V} \otimes \mathrm{V}^*\) 与 \(\operatorname{End}(\mathrm{V})\) 等同起来。证明存在 \(x \in \mathrm{V}, y \in \mathrm{V}^*\)，使得 \(x \otimes y \in \mathfrak{g}\) 且 \(\langle x, e^* \rangle \neq 0, \langle e, y \rangle \neq 0\)（取 \(\theta\) 在 \(e^n\) 下的共轭，其中 \(n\) 是 \(\mathfrak{g}\) 的一个适当的幂零元）。利用 \(\mathfrak{g}\) 是 \(\mathfrak{h}\)-子模（\(\mathrm{V} \otimes \mathrm{V}^*\) 中的）这一事实，推出 \(\mathfrak{g}\) 含有 \(e \otimes e^*\)。由此推出 \(\lambda + \lambda^{*} = \tilde{\alpha}\)，其中 \(\tilde{\alpha}\) 是 \(\mathfrak{s}\) 的最高根。
\end{enumerate}

\(d\) ) 证明 \(\mathfrak{s}\) 不是 \(\mathrm{B}_l\) 型（\(l\geq 2\)）、\(\mathrm{D}_l\) 型（\(l\geq 4\)）、\(\mathrm{E}_6,\mathrm{E}_7,\mathrm{E}_8,\mathrm{F}_4,\mathrm{G}_2\) 型（因为由第VI章的表，\(\tilde{\alpha}\) 将是基本权，从而不可能具有上面 \(\lambda +\lambda^{*}\) 的形式）。由此推出 \(\mathfrak{s}\) 是 \(\mathrm{A}_l\) 型或 \(\mathrm{C}_l\) 型，并且在第一种情形 \(\lambda = \varpi_1\) 或 \(\varpi_1 = \varpi_1^*\)，在第二种情形 \(\tilde{\alpha} = 2\varpi_1 = \varpi_1 + \varpi_1^*\)；由于 \(\mathfrak{c} = 0\) 或 \(k.1\)，这确实给出 \(a)^{21}\) 中的四种可能。

¶16) 令 g 为 \(A_{l}\) 型（\(l \geq 2\)）的绝对单李代数，\(\bar{k}\) 为 k 的代数闭包，\(\pi : \operatorname{Gal}(\bar{k}/k) \to \operatorname{Aut}(\bar{\mathbf{R}},\bar{\mathbf{B}})\) 为 §5 习题8 中定义的同态。

\begin{enumerate}
\def\labelenumi{\alph{enumi})}
\tightlist
\item
  假设 \(\pi\) 平凡。证明存在恰好两个双边理想 \(\mathfrak{m}\) 与 \(\mathfrak{m}'\)（\(\mathrm{U}(\mathfrak{g})\) 的），使得 \(\mathrm{D} = \mathrm{U}(\mathfrak{g}) / \mathfrak{m}\) 与 \(\mathrm{D}' = \mathrm{U}(\mathfrak{g}) / \mathfrak{m}'\) 都是维数为 \((l + 1)^2\) 的中心单代数（利用 §7 习题8）。\(\mathrm{U}(\mathfrak{g})\) 的主反自同构互换 \(\mathfrak{m}\) 与 \(\mathfrak{m}'\)；特别地，\(\mathrm{D}'\) 同构于 D 的反代数。复合映射 \(\mathfrak{g} \to \mathrm{U}(\mathfrak{g}) \to \mathrm{D}\) 把 \(\mathfrak{g}\) 与 D 中由迹为零的元素组成的李子代数 \(\mathfrak{sl}_{\mathrm{D}}\) 等同起来。此外，\(\mathfrak{g}\) 可分裂（从而同构于 \(\mathfrak{sl}(l + 1,k)\)）当且仅当 D 同构于 \(\mathbf{M}_{l + 1}(k)\)。
\end{enumerate}

反之，若 \(\Delta\) 是维数为 \((l + 1)^2\) 的中心单代数，则李代数 \(\mathfrak{sl}_{\Delta}\) 是 \(A_l\) 型的绝对单李代数，且相应的同态 \(\pi\) 平凡。两个这样的代数 \(\mathfrak{sl}_{\Delta}\) 与 \(\mathfrak{sl}_{\Delta'}\) 同构当且仅当 \(\Delta\) 与 \(\Delta'\) 同构或反同构。

\begin{enumerate}
\def\labelenumi{\alph{enumi})}
\setcounter{enumi}{1}
\tightlist
\item
  假设 \(\pi\) 非平凡。由于 \(\operatorname{Aut}(\bar{\mathbf{R}},\bar{\mathbf{B}})\) 有两个元素，\(\pi\) 的核是 \(\operatorname{Gal}(\bar{k}/k)\) 的指标为 2 的开子群，按 Galois 理论它对应于 k 的一个二次扩张 \(k_{1}\)。把 \(k_{1}\) 的非平凡对合记作 \(x\mapsto\bar{x}\)。
\end{enumerate}

证明存在唯一的双边理想 \(\mathfrak{m}\)（\(\mathrm{U}(\mathfrak{g})\) 的），使得 \(\mathrm{D} = \mathrm{U}(\mathfrak{g}) / \mathfrak{m}\) 是维数为 \(2(l + 1)^2\) 的单代数，且其中心是 \(k\) 的二次扩张（方法同上）。D 的中心可与 \(k_{1}\) 等同起来。理想 \(\mathfrak{m}\) 在 \(\mathrm{U}(\mathfrak{g})\) 的主反自同构下稳定；于是过渡到商即定义出 D 的一个对合反自同构 \(\sigma\)，使得 \(\sigma(x) = \bar{x}\) 对所有 \(x \in k_{1}\) 成立。复合映射 \(\mathfrak{g} \to \mathrm{U}(\mathfrak{g}) \to \mathrm{D}\) 把 \(\mathfrak{g}\) 与李子代数 \(\mathfrak{s}\mathfrak{u}_{\mathrm{D},\sigma}\)——由元素 \(x\)（满足 \(\sigma(x) = -x\) 且 \(\mathrm{Tr}_{\mathrm{D}/k_1}(x) = 0\)）组成——等同起来。反之，若 \(\Delta\) 是中心单 \(k_{1}\)-代数（维数为 \((l + 1)^2\)），且带有一个对合反自同构 \(\sigma\)（它在 \(k_{1}\) 上的限制为 \(x \mapsto \bar{x}\)），则李代数 \(\mathfrak{s}\mathfrak{u}_{\Delta,\sigma}\) 是 \(\mathrm{A}_l\) 型的绝对单李代数，且相应的同态 \(\pi\) 是与 \(k_{1}\) 相伴的那一个。两个这样的代数 \(\mathfrak{s}\mathfrak{u}_{\Delta,\sigma}\) 与 \(\mathfrak{s}\mathfrak{u}_{\Delta',\sigma'}\) 同构当且仅当存在 \(k\)-同构 \(f: \Delta \to \Delta'\) 使得 \(\sigma' \circ f = f \circ \sigma\)。

证明：当 \(\mathrm{D} = \mathbf{M}_{l+1}(k_{1})\) 时，存在 \(l + 1\) 阶的可逆埃尔米特矩阵 H，在相差一个 \(k^{*}\) 中元素乘积的意义下唯一，使得

\[
\sigma (x) = H. ^ {t} \bar {x}. H ^ {- 1}
\]

对所有 \(x \in \mathbf{M}_{l+1}(k_1)\) 成立；于是代数 g 可与代数 \(\mathfrak{su}(l + 1, H)\)（由满足 \(x.H + H.^t\bar{x} = 0\) 且 \(\operatorname{Tr}(x) = 0\) 的矩阵 x 组成）等同起来。

¶ 17) 令 g 为 \(B_{l}\) 型（分别地，\(C_{l}, D_{l}\) 型）的绝对单李代数，其中 \(l \geq 2\)（分别地，\(l \geq 3, l \geq 4\)）。进一步假设：当 g 为 \(D_{4}\) 型时，§5 习题8 中定义的同态 \(\pi : \operatorname{Gal}(\bar{k}/k) \to \operatorname{Aut}(\bar{\mathrm{R}}, \bar{\mathrm{B}})\) 的像的阶 \(\leq 2\)。证明存在维数为 \((2l + 1)^{2}\)（分别地，\(4l^{2}, 4l^{2}\)）的中心单代数 D 及 D 的一个对合反自同构 \(\sigma\)，使得 g 同构于 D 中由满足 \(\sigma(x) = -x\) 且 \(\operatorname{Tr}_{\mathrm{D}}(x) = 0\) 的元素 x 组成的李子代数（方法同习题16 a)）。\(^{22}\)

\begin{enumerate}
\def\labelenumi{\arabic{enumi})}
\setcounter{enumi}{17}
\tightlist
\item
  令 V 为有限维向量空间，Q 为 V 上的非退化二次型，\(\Psi\) 为与 Q 相伴的对称双线性形式。以 g 表示李代数 \(\mathfrak{o}(\varPsi)\)，并把 \(\varPsi\) 到 \(\bigwedge^{2}(V)\) 上的扩张记作 \(\varPsi_{2}\)。
\end{enumerate}

\begin{enumerate}
\def\labelenumi{\alph{enumi})}
\tightlist
\item
  证明存在向量空间同构 \(\theta : \bigwedge^{2}(\mathrm{V}) \to \mathfrak{g}\)，它由下列等价性质刻画：
\end{enumerate}

\begin{enumerate}
\def\labelenumi{(\roman{enumi})}
\item
  对 \(a, b\) 与 \(x\)（在 \(\mathrm{V}\) 中），有 \(\theta(a \wedge b). x = a. \Psi(x, b) - b. \Psi(x, a)\)。
\item
  对 \(x, y\)（在 V 中）及 \(u\)（在 \(\bigwedge^2 (\mathrm{V})\) 中），有 \(\varPsi_2(x\wedge y,u) = \varPsi (x,\theta (u).y)\)。
\end{enumerate}

令 \(\sigma\) 为 \(\mathfrak{g}\) 在 \(V\) 上的恒等表示；则 \(\theta\) 是 \(\bigwedge^{2}(V)\) 与 \(\mathfrak{g}\) 的伴随表示之间的同构。

\begin{enumerate}
\def\labelenumi{\alph{enumi})}
\setcounter{enumi}{1}
\item
  按第2节引理1 的方式定义线性映射 \(f: \mathfrak{g} \to \mathrm{C}^{+}(\mathrm{Q})\)。证明对 V 中的 a, b 有 \(f\theta(a \wedge b) = \frac{1}{2}(ab - ba)\)，并由此给出第2节引理1 的断言 (iii)、(iv) 与 (v) 的一个新证明。
\item
  记号 \(l, \mathrm{F}, \mathrm{F}', e_0, \mathrm{N}, \lambda\) 与 \(\rho\) 沿用第2节的。取 \(e \neq 0\)（在 \(\bigwedge (\mathrm{F}')\) 中），并用下式定义 \(\varPhi\)（\(\mathrm{N}\) 上的双线性形式）：
\end{enumerate}

\[
x \wedge y = (- 1) ^ {p (p + 1) / 2} \Phi (x, y). e \quad (x \in \bigwedge^ {p} (\mathrm{F} ^ {\prime}), y \in \mathrm{N}).
\]

证明 \(\Phi(\lambda(a).x,y)+\Phi(x,\lambda(a).y)=0\)（对 N 中的 x, y 及 \(F\oplus F'\) 中的 a）。由此推出 \(\Phi\) 在 g 的表示 \(\rho\) 下不变（注意由上面的 a) 与 b)，李代数 \(\rho(\mathfrak{g})\) 由 \(\lambda(F\oplus F')\) 生成）。

\begin{enumerate}
\def\labelenumi{\arabic{enumi})}
\setcounter{enumi}{18}
\tightlist
\item
  令 \(\mathrm{V}\) 为有限维向量空间，\(\mathrm{E} = \mathrm{V} \oplus \mathrm{V}^*\)。用下式定义 \(\varPhi\)（\(\mathrm{E}\) 上的非退化双线性形式）：
\end{enumerate}

\[
\Phi ((x, x ^ {*}), (y, y ^ {*})) = \langle x, y ^ {*} \rangle + \langle y, x ^ {*} \rangle .
\]

置 \(N = \bigwedge(V)\)，并置 \(Q(x) = \frac{1}{2}\Phi(x, x)\)（对 \(x \in E\)）。以 \(\lambda : C(Q) \to \text{End}(N)\) 表示第2节.IV 中的旋量表示，按第2节引理1 定义 \(f : \mathfrak{o}(\Phi) \to C^{+}(Q)\)，并以 \(\rho\) 表示线性表示 \(\lambda \circ f\)（代数 \(\mathfrak{o}(\Phi)\) 在 N 上的）。

\begin{enumerate}
\def\labelenumi{\alph{enumi})}
\item
  对 \(u\)（\(V\) 的任一自同态），用 \(\tilde{u}\)（\(E\) 的自同态）与之对应，公式为 \(\tilde{u}(x, x^{*}) = (u(x), -^{t}u(x^{*}))\)。证明 \(u \mapsto \tilde{u}\) 是从 \(\mathfrak{gl}(V)\) 到 \(\mathfrak{o}(\varPhi)\) 的李代数同态；此外，对 \(u\)（\(\mathfrak{gl}(V)\) 中的），\(\rho(\tilde{u})\) 是代数 \(\bigwedge(V) = N\) 的唯一导子，它与 \(u\)（在 \(V\) 上）一致。
\item
  令 \(\varPsi\) 为 V 上的非退化交错双线性形式，\(\gamma: V \to V^{*}\) 为由 \(\varPsi(x, y) = \langle x, \gamma(y) \rangle\)（x, y 在 V 中）定义的同构。证明由下式定义的 E 的自同态 \(\tilde{X}_{+}\) 与 \(\tilde{X}_{-}\)
\end{enumerate}

\[
\tilde {X} _ {+} (x, x ^ {*}) = (\gamma^ {- 1} (x ^ {*}), 0), \quad \tilde {X} _ {-} (x, x ^ {*}) = (0, - \gamma (x))
\]

属于 \(\mathfrak{o}(\varPhi)\)。置 \(\tilde{H} = (-1)^{\sim}\)。证明 \((\tilde{H}, \tilde{X}_{+}, \tilde{X}_{-})\) 是一个 \(\mathfrak{sl}_2\) 三元组（在李代数 \(\mathfrak{o}(\varPhi)\) 中）。

\begin{enumerate}
\def\labelenumi{\alph{enumi})}
\setcounter{enumi}{2}
\tightlist
\item
  证明 \(\rho\) 把自同态 \(\tilde{H}, \tilde{X}_{+}, \tilde{X}_{-}\)（\(\mathfrak{o}(\Phi)\) 的）分别变为第3节.IV 中记作 \(H, X_{+}\) 与 \(X_{-}\) 的 N 的自同态。由此推出 \((H, X_{+}, X_{-})\) 是一个 \(\mathfrak{sl}_{2}\) 三元组（在李代数 \(\mathfrak{gl}(N)\) 中）。
\end{enumerate}

